\chapter{Chemotherapy}

\section*{Overview}
Chemotherapy uses chemical substances (antineoplastic drugs) to kill rapidly dividing cancer cells. It can be systemic, affecting cells throughout the body. The exact approach, setting, and preparation depend on the indication, anatomy, and the patient's health.

\section*{Alternative Names}
Chemo, Cancer chemotherapy

\section*{Principle}
Cytotoxic drugs are given orally or intravenously and circulate through the body, where they act on rapidly dividing cells by damaging DNA, blocking cell division, or interrupting signaling pathways. The goal is to kill cancer cells while normal tissues recover between treatment cycles.
\section*{History}
Began in the 1940s with the use of nitrogen mustard (derived from chemical warfare research) to treat lymphoma. Sidney Farber pioneered folic acid antagonists for leukemia.

\section*{Why It Is Done}
Ordered to cure cancer, control cancer growth, shrink tumors before surgery (neoadjuvant), or kill remaining cancer cells after surgery (adjuvant).

\section*{Is This a Primary or Secondary Test or Procedure}
Primary treatment for hematological malignancies and advanced solid tumors.

\section*{How It Is Performed}
Administered in cycles, usually intravenously through a central line or port-a-cath, or orally as pills. Treatment is given over several hours or days, followed by a rest period.

\section*{Preparation}
Baseline blood tests (CBC, liver/kidney function) and cardiac evaluation (echocardiogram) are required before each cycle. Hydration is key.

\section*{Contraindications and Precautions}
Severe bone marrow suppression (absolute neutrophil count <1,000 or platelets <50,000), severe active infection, or end-organ dysfunction depending on the drug.

\section*{Risks}
Neutropenia (infection risk), thrombocytopenia (bleeding), anemia, nausea, hair loss, fatigue, cardiotoxicity, or peripheral neuropathy.

\section*{Aftercare and Recovery}
Monitor for fever (>100.4 F) which indicates neutropenic fever, an oncology emergency. Manage side effects with antiemetics and supportive care.

\section*{Outcome and Follow-up}
Monitored using imaging scans (CT, MRI) and tumor markers to assess tumor shrinkage or remission.

\section*{Expected Results}
Stable or shrinking tumor size; normal bone marrow recovery between cycles.

\section*{Follow-up}
CBC monitoring, imaging scans, and regular oncology follow-up.

\section*{When to Seek Medical Care}
Follow the procedure team's specific instructions. Seek urgent medical assessment for severe or worsening pain, uncontrolled bleeding, fever or signs of infection, trouble breathing, chest pain, fainting, new weakness or confusion, inability to urinate, or any rapidly worsening symptom. The appropriate warning signs depend on the procedure and the patient's medical condition.

\section*{References}
\begin{enumerate}
  \item MedlinePlus. Chemotherapy. U.S. National Library of Medicine; 2025.
  \item Current Medical Diagnosis and Treatment. McGraw-Hill; 2025.
\end{enumerate}

\clearpage
